\documentclass{book}

\title{COLA}
\newcommand{\booksubtitle}{computational oriented linear algebra}
\newcommand{\booklicense}{ALL RIGHT PRESERVED}

\author{Zhang Jinrui}
% Author subtitle could be a university or a geographical location, for example
\newcommand{\authorsubtitle}{NewArea, China}

% Create convenient commands \booktitle and \bookauthor
\makeatletter
\newcommand{\booktitle}{\@title}
\newcommand{\bookauthor}{\@author}
\makeatother

% This utf8 declaration is not needed for versions of latex > 2018 but may
% be helpful for older software. Eventually it may not be worth keeping.
\usepackage[utf8]{inputenc}  
\usepackage{fix-cm}  % this package allows large \fontsize
\usepackage{tikz}    % this is for graphics. e.g. rectangle on title page
\usepackage{amsmath} % Used by equations
\usepackage{amssymb}
% The following dimensions specify 4.75" X 7.5" content on 6 3/8" by 9 1/4"
% paper. The paper width and height can be tweaked as required and the content
% should size to fit within the margins accordingly.
%
% The (inside) bindingoffset should be larger for books with more pages. Some
% standard recommended sizes are .375in minimum up to 1in for 600+ page books.
% Sizes .75in and .875in are also recommended roughly at 150 and 400 pages.
\usepackage[bindingoffset=0.625in,
            left=.5in, right=.5in,
            top=.8125in, bottom=.9375in,
            paperwidth=6.375in, paperheight=9.25in]{geometry}
% Here is an alternative geometry for reading on letter size paper:
% \usepackage[margin=.75in, paperwidth=8.5in, paperheight=11in]{geometry}

\renewcommand{\contentsname}{Table of Contents} % default is {Contents}
\usepackage{makeidx}
\makeindex % Initialize an index so we can add entries with \index

% The next few commands are for creating fake content to fill out the template.
% You should delete this (e.g.  everything up to, but not including,
% \begin{document}) after you insert your own content.
% Example content from Einstein's Meaning of Relativity.
% Public domain book: http://www.gutenberg.org/ebooks/36276
\newcommand{\fakeparagraph}{}
\newcommand{\para}{defsdf}
\newcommand{\fakecontent}{
    sldkfjsldkfj

    sdfsdfsdf
}
\newcommand{\prefacePara}
{
The current 
textbook in linear algebra is very abundant, 
include the linear operator oriented approach, 
which is the currently best one I have ever seen.

This book will base on a computational oriented 
approach to the subject of the so called linear
algebra.Linear algebra is such a vast topic that
I'm having the most firghtened heart to output
what I have learnt so far in this amazing land.

In this book, the main goal is to answer all the
questions about those are not well pondered in the
traditional linear algebra course of any kind.
In the computational matrix theory there are plenty
ways of decompose the matrix, such as LU, QR, and the
most fundamental SVD decomposition.
While in the more math oriented linear course
the main topic are always in the most abstract way
as they move forward.
And the main concepts such as the Transpose, dule
space, determinant, trace and all other fundamental
but always treated just in a forced memorize level.

The main goal is to explain all the concepts in a 
more mathematical natrual and fun way, even the sudden
expose of a certain defination would be a very hard
problem to those first expose to the abstract algebra
world student just like me. So I'll try my best to
get every defination and notion in a smooth and more
reasonalbe way. 
}

\newcommand{\fakesections}{
\fakecontent{}
\section{Tensor oriented for computation}\fakecontent{}
\subsection{Examples}\fakecontent{}
\subsection{Excessive Elaborations}\fakecontent{}
\subsection{Long Winded Conclusion}\fakecontent{}
\subsection{Exercises}\fakecontent{}
\section{Theory vs Practice}\fakecontent{}
\subsection{Examples}\fakecontent{}
\subsection{Excessive Elaborations}\fakecontent{}
\subsection{Long Winded Conclusion}\fakecontent{}
\subsection{Exercises}\fakecontent{}
\foreach \x in {A,B,C,D,E,F,G,H,I,J,K,L,M,N,O,P,Q,R,S,T,U,V,W,X,Y,Z}
    {\index{\x 1}\index{\x 2}}
}

% Content Starts Here
\begin{document}
\frontmatter

% ---- Half Title Page ----
% current geometry will be restored after title page
\newgeometry{top=1.75in,bottom=.5in}
\begin{titlepage}
    \begin{flushleft}

        % Title
        \textbf{\fontfamily{qcs}\fontsize{48}{54}\selectfont COLA\\}

        % Draw a line 4pt high
        \par\noindent\rule{\textwidth}{4pt}\\

        % Subtitle
        % Shaded box from left to right. Text node is midway (centered).
        \begin{tikzpicture}
            \shade[bottom color=lightgray,top color=white]
            (0,0) rectangle (\textwidth, 1.5)
            node[midway] {\textbf{\large \textit{\booksubtitle}}};
        \end{tikzpicture}

        % Edition Number
        \begin{flushright}
            \Large First Edition
        \end{flushright}

        % \vspace{\fill}
        \vspace{\fill}

    \end{flushleft}
\end{titlepage}
\restoregeometry
% ---- End of Half Title Page ----


% No page numbers on the Frontispiece page
\thispagestyle{empty}


% ---- Title Page ----
% current geometry will be restored after title page
\newgeometry{top=1.75in,bottom=.5in}
\begin{titlepage}
    \begin{flushleft}

        % Title
        \textbf{\fontfamily{qcs}\fontsize{48}{54}\selectfont COLA\\}

        % Draw a line 4pt high
        \par\noindent\rule{\textwidth}{4pt}\\

        % Shaded box from left to right with Subtitle
        % The text node is midway (centered).
        \begin{tikzpicture}
            \shade[bottom color=lightgray,top color=white]
            (0,0) rectangle (\textwidth, 1.5)
            node[midway] {\textbf{\large \textit{\booksubtitle}}};
        \end{tikzpicture}

        % Edition Number
        \begin{flushright}
            \Large First Edition
        \end{flushright}

        \vspace{\fill}

        % Author and Location
        \textbf{\large \bookauthor}\\[3.5pt]
        \textbf{\large \textit{\authorsubtitle}}

        \vspace{\fill}

        % Self Publishing Logo. Free to use: CC0 license.
        % The source file is book.svg. If you change the svg, you must then convert
        % it to pdf. There are many online and offline tools available to do that.
        \begin{center}
            \includegraphics{booksvg.pdf}\\[4pt]
            \fontfamily{lmtt}\small{Self Publishers Worldwide\\
                Seattle San Francisco New York\\
                London Paris Rome Beijing Barcelona}
        \end{center}

    \end{flushleft}
\end{titlepage}
\restoregeometry
% ---- End of Title Page ----

% Do not show page numbers on colophon page
\thispagestyle{empty}

\begin{flushleft}
    \vspace*{\fill}
    This book was typeset using \LaTeX{} software.\\
    To those who love me.\\
    \vspace{\fill}
    Copyright \textcopyright{} \the\year{}  \bookauthor\\
    License: \booklicense
\end{flushleft}

% A title page resets the page # to 1, but the second title page
% was actually page 3. So add two to page counter.
\addtocounter{page}{2}

% The asterisk excludes chapter from the table of contents.
\chapter*{Preface}
\prefacePara{}
% \[
%     \left.
%     \begin{aligned}
%          & \sqrt{{dX_{1}}^{2} + {dX_{2}}^{2} + {dX_{3}}^{2}} \\
%          & \qquad
%         \begin{aligned}
%                & = \left(1 + \frac{\kappa}{8\pi} \int \frac{\sigma\, dV_{0}}{r}\right)
%             \sqrt{{dx_{1}}^{2} + {dx_{2}}^{2} + {dx_{3}}^{2}},                             \\
%             dT & = \left(1 - \frac{\kappa}{8\pi} \int \frac{\sigma\, dV_{0}}{r}\right) dl.
%         \end{aligned}
%     \end{aligned}
%     \right\}
% \]

% Three-level Table of Contents
\setcounter{tocdepth}{3}
\tableofcontents

\mainmatter

\chapter{PreTalks}
 {
  the first thing to think about is how we mearsure things.
  \section{Tensor oriented for computation}
  \subsection{Examples of (a,b)-Tensor $\forall a,b \in \mathbb{N}$}
  \subsection{(0,1)-Tensor the vector}
  \subsection{(1,0)-Tensor the covector}
  \subsection{(1,1)-Tensor the linear maps}
  \subsection{(2,0)-Tensor the bilinear maps}
  \subsection{Exercises}
  \section{Outer algebra oriented for computation}
  \subsection{Examples of geometric algebra}
  \subsection{bivectors and trivectors}
  \subsection{dot and wedge product and the flux in 3D}
  \subsection{Exercises}
  \foreach \x in {A,B,C,D,E,F,G,H,I,J,K,L,M,N,O,P,Q,R,S,T,U,V,W,X,Y,Z}
      {\index{\x 1}\index{\x 2}}
 }
\chapter{Tensor(Vector) Space}
 {
  the first thing to think about is how we mearsure things.
  \section{Tensor oriented for computation}
  \subsection{Examples of (a,b)-Tensor $\forall a,b \in \mathbb{N}$}
  \subsection{(0,1)-Tensor the vector}
  \subsection{(1,0)-Tensor the covector}
  \subsection{(1,1)-Tensor the linear maps}
  \subsection{(2,0)-Tensor the bilinear maps}
  \subsection{Exercises}
  \section{Outer algebra oriented for computation}
  \subsection{Examples of geometric algebra}
  \subsection{bivectors and trivectors}
  \subsection{dot and wedge product and the flux in 3D}
  \subsection{Exercises}
  \foreach \x in {A,B,C,D,E,F,G,H,I,J,K,L,M,N,O,P,Q,R,S,T,U,V,W,X,Y,Z}
      {\index{\x 1}\index{\x 2}}
 }
\chapter{Irrational and Transcendent Numbers}
\chapter{Complex Numbers}
\chapter{Quaternions and Ausdehnungslehre}
\chapter{Theory of Equations}
\chapter{Substitutions and Groups}
\chapter{Determinants}
\chapter{Quantics}
\chapter{Calculus}
\chapter{Differential Equations}
\chapter{Infinite Series}
\chapter{Theory of Functions}
\chapter{Probabilities and Least Squares}
\chapter{Analytic Geometry}
\chapter{Modern Geometry}
\chapter{Trigonometry and Elementary Geometry}
\chapter{Non-Euclidean Geometry}
\chapter{Bibliography}
\chapter{General Tendencies}

\backmatter
\addcontentsline{toc}{chapter}{Index}
\printindex
\end{document}
